\chapter{Differences and Known Limitations}
\label{app:differences}

Everything \ccprod{} does not do, or does differently, in one list. Two kinds of
entry are mixed here and marked apart: \textbf{absences}, which are work that
has not been done, and \textbf{defects}, which are work that was done wrong.

\section{Absences}

\subsection{Selection and editing}

\begin{ccsteps}
\item \textbf{No fill, no fill down, no fill series.} A column of similar
formulas is built by copying a block and pasting it, which does adjust the
references --- but there is no one-key fill.
\item \textbf{No cut or move.} Copy and paste, then clear what you left.
\item \textbf{No multi-level undo.} One cell, one level, and a pasted or
cleared block cannot be undone at all.
\item \textbf{No Find Next.} There is no need for one: repeating
\menupath{Find} walks the matches. See Chapter~\ref{ch:finding}.
\item \textbf{Replace does not touch formulas}, and cannot be undone.
\item \textbf{No Go to.} You cannot type an address to jump to it.
\item \textbf{No insert or delete of a single cell.} Whole rows and whole
columns only; neighbours cannot be shifted on their own.
\end{ccsteps}

\subsection{Formatting}

\begin{ccsteps}
\item \textbf{No formatting commands of any kind.} Number formats, bold,
alignment and column widths are all supported by the display and both file
formats --- and none of them can be set from the keyboard. They arrive from
imported files only.
\item \textbf{One currency symbol, \lit{\$}}, and no thousands separators.
\item \textbf{One date format, \lit{YYYY-MM-DD}.}
\item \textbf{No colours, borders, fonts or cell shading.}
\item \textbf{No split screen, no zoom, no hidden rows or columns.} There
\emph{are} freeze panes --- see Chapter~\ref{ch:moving}.
\item \textbf{No printing.}
\end{ccsteps}

\subsection{Formulas}

\begin{ccsteps}
\item \textbf{Fourteen functions.} No text functions, lookup functions,
statistical functions beyond \fn{AVERAGE}, financial functions or date
functions. Chapter~\ref{ch:funcref} has the full list of what is missing.
\item \textbf{No \lit{\&} concatenation} and no \lit{\%} postfix.
\item \textbf{No named ranges}, no whole-column references (\lit{A:A}), no
three-dimensional ranges across worksheets, no array formulas.
\item \textbf{No R1C1 notation.}
\item \textbf{No iterative calculation.} (There \emph{is} a manual
recalculation mode --- see Chapter~\ref{ch:recalc}.)
\item \textbf{Semicolon is not accepted} as an argument separator, and a tab is
not treated as whitespace.
\end{ccsteps}

\subsection{Elsewhere}

\begin{ccsteps}
\item \textbf{No macros, windows, or linked workbooks.}
\item \textbf{Three chart types}, all of one column, at most sixteen values:
bars, a line and a pie. None is stored in the workbook, though any of them can
be written out as a \fname{.BMP}. No stacked, scatter or area charts.
Chapter~\ref{ch:charting}.
\item \textbf{No way to change device or directory} from inside the program.
\item \textbf{No way to type a date.}
\end{ccsteps}

\section{Differences in behaviour}

Things that exist here and work differently from the way another spreadsheet
would do them.

\begin{cctablethree}{Case}{Commander Calc}{Elsewhere}{1.5in}{1.55in}{1.77in}
\fml{=2\^{}3\^{}2} & 64 & 512 in mathematics; 64 in Excel \\
\fml{=-2\^{}2} & 4 & $-4$ in mathematics; 4 in Excel \\
\fn{MIN} or \fn{MAX} of nothing & 0 & 0 in Excel \\
\fn{AVERAGE} of nothing & \errv{\#DIV/0!} & \errv{\#DIV/0!} in Excel \\
\fn{MOD} sign & Takes the divisor's & Same in Excel; differs from C \\
\fn{INT} of $-2.7$ & $-3$ & $-3$ in Excel; $-2$ if truncating \\
Circular reference & \errv{\#CYCLE!} in the cell & 0 and a status warning in
Excel \\
A number below 1 & \lit{.5}, no leading zero & \lit{0.5} \\
Value too wide for a column & Cut off silently & \lit{\#\#\#\#} \\
\lit{\$} in a reference & Parsed and ignored & Fixes the reference on copy \\
Deleting a referenced row & Points at the replacement & \errv{\#REF!} \\
\end{cctablethree}

Where \ccprod{} and Excel agree, that is deliberate: the point of an import
format is that the numbers come out the same.

\section{Defects}

These are wrong rather than absent. They are documented here because a manual
that quietly omitted them would be worse than useless.

\subsection{END goes to the wrong place}

\key{End} moves to column \cellref{IV} --- the last column that can exist ---
rather than to the last column in use. In an empty worksheet that is a jump of
255 columns into nothing. \keyc{Shift}{Home} gets you back.

\subsection{A sorted range is not rewritten, and cannot be}

Inserting rows, deleting rows, inserting and deleting columns, sorting, undoing
a sort and pasting all rewrite the references inside formulas so that they go on
meaning the same cells. One case is beyond that treatment.

\begin{cccaution}
\textbf{A range spanning rows that a sort moves is wrong afterwards.}
\fml{=SUM(B2:B10)} names nine consecutive rows; a sort scatters those rows, and
no rewriting can express where they went, because they are no longer next to
each other. \ccprod{} asks before sorting rows that hold formulas ---
\msg{ranges over sorted rows may be wrong} --- and that is the case it is
asking about.

Keep range formulas out of the rows being sorted: in a summary block above the
table, or on another worksheet.
\end{cccaution}

Rewriting has three further gaps, all documented in Chapter~\ref{ch:rows}: a
reference naming another worksheet is left alone; a reference pointing straight
at a deleted row silently ends up on whatever replaced it, where Excel would
show \errv{\#REF!}; and a formula over 128 characters long is not rewritten at
all.

\subsection{Text comparison is by identity, not by content}

\begin{cccaution}
Two cells both containing \lit{abc} compare as unequal, because each piece of
text in a workbook is a separate object and comparison tests whether they are
the same object. \fml{="a"="a"} is \msg{FALSE}.

Worse, unequal text compares as greater in both directions:
\fml{="a">"b"} and \fml{="b">"a"} are both \msg{TRUE}, and \lit{<} on text is
always \msg{FALSE}.

Comparison is reliable on numbers, dates and booleans only.
\end{cccaution}

\subsection{The time format is not implemented}

A cell formatted as a time displays its raw fractional number rather than a
time. This is deliberate --- visibly unfinished was preferred to plausibly
wrong --- but it is a gap, not a design.

\subsection{A date-formatted cell holding an out-of-range value}

A cell with a date format whose value is negative, or too large to be a day
number, produces unpredictable rubbish rather than an error. Day number 0
displays as \lit{1900-01-00}, which is at least recognisable as nonsense.

\subsection{The note on an import result box is wrong}

The counts on an import result box are correct. The parenthesised note after
them is off by one, so the wrong phrase can appear --- or an empty pair of
brackets. \textbf{Read the counts, not the note.}

\subsection{Cosmetic}

\begin{ccsteps}
\item The highlight on an open menu's name is four characters wide whatever the
name is, so \menupath{Sheet} shows as \msg{Shee} with a stray \msg{t}.
\item One character of the formula bar, between the separator and the contents,
is never painted and shows through in the worksheet colour.
\item Mouse hit-testing inside menus and dialogue boxes assumes an 80 by 60
display, and is inaccurate in any other text mode.
\item The Import file list and its busy box are both titled \msg{Open}.
\item A failure in \menupath{New} or \menupath{Save as} is reported as
\msg{Save failed}.
\end{ccsteps}

\section{Limits that will surprise you}

Not defects, but sharp edges worth knowing about in advance.

\begin{cctable}{Limit}{Consequence}{1.65in}{3.17in}
32 files in a dialogue & A directory with more workbooks than that has some
that cannot be opened, with no indication. \\
95 characters per CSV field & Longer text is truncated silently on import. \\
80 characters per entry & Long formulas must be built in a file and imported. \\
10 live values per formula & A function with eleven or more arguments gives
\errv{\#ERR11!}. \\
128K per XLSX part & Bigger than the cell limit suggests, and it is about the
file rather than the data. \\
512 formulas per XLSX import & The rest keep their cached values. \\
4,096 rows per sort & \msg{too many rows to sort}. \\
15 characters at the Find prompt & \\
\end{cctable}
